% Generated from locale/tr/content/lambda-calculus/church-rosser/beta-eta-reduction.tex; do not hand-edit.


\olfileid[tr]{lam}{cr}{be}

\olsection{$\beta\eta$-indirgeme}

Church--Rosser özelliği $\beta\eta$-indirgeme ($\bered$) için geçerlidir.

\begin{lem}\ollabel{lem:one-par}
  $M \beredone M'$ ise $M \beredpar M'$.
\end{lem}

\begin{proof}
  $M \beredone M'$ !!{derivation} üzerinde tümevarımla. Eğer
  $M \eredone M'$, yani bu adım $\eta$-dönüşümüyle
  (\olref[syn][eta]{defn:beredone}) elde edilmişse
  \olref[pbe]{thm:refl} kullanılır. Diğer durumlar
  \olref[b]{lem:one-par} içindeki gibidir.
\end{proof}


\begin{lem}\ollabel{lem:par-red}
  $M \beredpar M'$ ise $M \bered M'$.
\end{lem}

\begin{proof} $M \beredpar M'$ !!{derivation} üzerinde tümevarımla.

  Son kural \olref[pbe]{defn:beredpar5} ise bazı $x,N,N'$ için
  $M=\lambd[x][Nx]$, $M'=N'$, $x\notin\FV{N}$ ve $N\beredpar N'$ olur.
  Önce $\lambd[x][Nx]$ ifadesini $\eta$-dönüşümüyle $N$'ye, ardından
  tümevarım varsayımına göre $N\bered N'$ bağıntısını veren $\beredone$
  adımları dizisiyle $N'$ye indirgeyebiliriz.
\end{proof}


\begin{lem}\ollabel{lem:str}
  $\bered$, $\beredpar$ bağıntısını içeren en küçük geçişli bağıntıdır.
\end{lem}

\begin{proof}
  \olref[b]{lem:str} içindeki gibi.
\end{proof}

\begin{thm}\ollabel{thm:cr}
  $\bered$ Church--Rosser özelliğini sağlar.
\end{thm}

\begin{proof}
  \olref[dap]{thm:str}, \olref[pbe]{thm:cr} ve \olref{lem:str} gereği.
\end{proof}
